% ===================================================================
%  TAULA N° 5 — L'Ostal e sa construccion.
%  Traduction 2026 d'apres texto/io, controlee sur texto/fr et
%  texto/ca. Consigne : voir texto/oc/10-taula-01.tex.
%
%  Le tableau 5 est un DIALOGUE : on garde \textsc{} pour le nom du
%  parleur. Les renvois du plan sont des LETTRES, (a) a (m), et l'on
%  ecrit ce que l'ido ecrit, « (1) » compris la ou l'imprimeur a
%  compose un 1 pour un l.
%
%  LES DEUX NOMS PROPRES SE TRADUISENT, ET L'IDO L'A DEJA FAIT :
%  Leblanc est « Blanko », Roux est « Rufo ». L'occitan a « Blanc » et
%  « Roge », l'un et l'autre patronymes courants du Midi.
% ===================================================================

%%K t05-apar-1 apar
\VUsaut{6.0mm}
\VUtitre{15.0pt}{60}{TAULA N° 5}
\VUsaut{1.4mm}
\VUtitre{11.4pt}{0}{L'Ostal e sa construccion.}
\VUsaut{2.2mm}

%%K t05-tit-1 sub
\VUpk{10.2pt}{40}{Una bona rescontrada.}
\VUsaut{3.0mm}

%%K t05-01-1 p
\textsc{Joan}. --- Oncle, m'agradariá fòrça de saber cossí se construís un
\VUgras{ostal}.

%%K t05-01-2 p
\textsc{L'Oncle} \textsuperscript{(80)}. --- Es fòrça aisit, filh meu; vèni amb
ieu e te mostrarai aquel que lo sénher Bernat \textsuperscript{(79)} fa
construire e ont demorarai.

%%K t05-01-3 p
\textsc{Joan}. --- E qual a dessenhat lo \VUgras{plan} \textsuperscript{(76)}
d'aqueste ostal?

%%K t05-01-4 p
\textsc{L'Oncle}. --- L'\VUgras{arquitècte} \textsuperscript{(75)}; presentèt un
dessenh fòrça clar, que foguèt aprovat, e l'\VUgras{entrepreneire}
\textsuperscript{(77)} comencèt los obratges.

%%K t05-01-5 p
\textsc{Joan}. --- E qual es l'entrepreneire?

%%K t05-01-6 p
\textsc{L'Oncle}. --- Es lo sénher Blanc, lo paire de ton amic Jaume. Dirigís los
obrièrs amb lo sénher Roge, l'arquitècte.

%%K t05-01-7 p
Agacha! Vaicí que ven justament lo sénher Blanc amb sa \VUgras{règla}
\textsuperscript{(78)}, e lo sénher Roge amb son plan.

%%K t05-01-8 p
Bonjorn, senhers. Cossí anatz?

%%K t05-01-9 p
\textsc{Lo Sr. Roge}. --- Fòrça plan, mercés. Bonjorn, Joan; qu'a creissut aquel
enfant!

%%K t05-01-10 p
\textsc{L'Oncle}. --- O cresi ben, es tot un omenet. Es fòrça seriós; li cal
explicar tot çò que vei. Uèi voldriá saber cossí se construís un ostal.

%%K t05-tit-2 sub
\VUpk{10.2pt}{40}{Los obrièrs.}
\VUsaut{3.0mm}

%%K t05-01-11 p
\textsc{Lo Sr. Blanc}. --- Alara vèni amb ieu, amiguet, e t'o explicarai
sulcòp, que m'agradan fòrça los enfants que vòlon aprene.

%%K t05-01-12 p
Veses aquel obrièr qu'a una \VUgras{pala} \textsuperscript{(82)} a la man: es un
\VUgras{terrassièr} \textsuperscript{(81)}. Es a dobrir una \VUgras{trencada}
\textsuperscript{(46)} per i pausar la conducha de l'aiga. Es tanben el que cavèt
lo sòl a l'endrech ont es l'ostal, per que lo maçon auce las quatre
\VUgras{parets}. La partida d'aquelas parets que demòra jos tèrra s'apèla los
\VUgras{fondaments} \textsuperscript{(2)}. L'espaci compres entre los fondaments
forma la \VUgras{cava} \textsuperscript{(3)}. Aquela cava a una fenestreta apelada
\VUgras{lucanèl} \textsuperscript{(5)}, una \VUgras{pòrta}
\textsuperscript{(4)} e un escalièr.

%%K t05-01-13 p
Lo \VUgras{maçon} \textsuperscript{(108)} contunhèt d'auçar las parets en daissant
de traucs per las \VUgras{pòrtas} \textsuperscript{(8)} e las \VUgras{fenèstras}
\textsuperscript{(17)}.

%%K t05-01-14 p
\textsc{Joan}. --- Mas lo vesi pas, aquel maçon!

%%K t05-01-15 p
\textsc{Lo Sr. Blanc}. --- A ja acabat de trabalhar a l'ostal. Es contra
l'\VUgras{estable} \textsuperscript{(54)}, sus son \VUgras{escafaud}
\textsuperscript{(110)}, a auçar la paret del \VUgras{lavador}
\textsuperscript{(56)}.

%%K t05-01-16 p
A una truèla, una gaveta e un \VUgras{plomb} \textsuperscript{(109)}. Mas te
sovendràs pas d'aqueles noms.

%%K t05-01-17 p
\textsc{Joan}. --- Que òc que me'n sovendrai, que demandarai a mon oncle una
truèla pichona e una gaveta, e ieu meteis farai un plomb amb un cordèl.

%%K t05-tit-3 sub
\VUsaut{3.0mm}
\VUpk{10.2pt}{40}{Los materials.}
\VUsaut{3.0mm}

%%K t05-01-18 p
\textsc{L'Oncle}. --- A, fòrça plan! Mas digas, Joan, qué cal per construire un
ostal?

%%K t05-01-19 p
\textsc{Joan}. --- Cal de \VUgras{pèiras talhadas} \textsuperscript{(59)} e de
\VUgras{mortièr} \textsuperscript{(65)}.

%%K t05-01-20 p
\textsc{L'Oncle}. --- Exactament. Agacha aquel òme al grand capèl, sus son
\VUgras{carri} \textsuperscript{(136)}. Es un \VUgras{carretièr}
\textsuperscript{(135)}. Cada jorn pòrta aicí de \VUgras{blòts de pèira}
\textsuperscript{(60)}. Descarga son carri dins la cort, e los
\VUgras{peirièrs} \textsuperscript{(83)} talhan los blòts e los rèsson amb sa
\VUgras{rèssa de pèira} \textsuperscript{(85)}. Un \VUgras{manòbra}
\textsuperscript{(146)} los pòrta al maçon, que los pausa.

%%K t05-01-21 p
Pendent aqueste temps, lo \VUgras{gachaire} \textsuperscript{(87)} prepara lo
mortièr. Li cal de \VUgras{sabla} \textsuperscript{(62)}, de \VUgras{cauç}
\textsuperscript{(63)} e d'\VUgras{aiga} \textsuperscript{(64)}. La cauç es a son
costat dins un \VUgras{sac} \textsuperscript{(71)}; un \VUgras{ajudaire de maçon}
\textsuperscript{(111)} li pòrta la sabla dins una \VUgras{bèrla}
\textsuperscript{(112)}, e l'\VUgras{aprentís} \textsuperscript{(113)} es a la
bomba \textsuperscript{(58)}. Emplena un \VUgras{selh} \textsuperscript{(114)}. Lo
mortièr serà lèu prèst. Lo \VUgras{portaire de mortièr} \textsuperscript{(107)}
lo pren e lo puja per l'escala fins a l'escafaud, ont un maçon acaba aquel costat
de l'ostal.

%%K t05-01-22 p
E ara, cossí s'apèla l'obrièr que fa lo \VUgras{teulat}
\textsuperscript{(31)}?

%%K t05-01-23 p
\textsc{Joan}. --- Es lo \VUgras{teulejaire} \textsuperscript{(118)}.

%%K t05-tit-4 sub
\VUsaut{3.0mm}
\VUpk{10.2pt}{40}{Lo teulat de l'ostal.}
\VUsaut{3.0mm}

%%K t05-01-24 p
\textsc{Lo Sr. Blanc}. --- Òc, mas d'en primièr ven lo \VUgras{fustièr de gròs}
\textsuperscript{(115)}.

%%K t05-01-25 p
Lo fustièr de gròs fa la \VUgras{cabirada} \textsuperscript{(35)}. Unís las
\VUgras{bigas} \textsuperscript{(37)} amb de \VUgras{cavilhas}
\textsuperscript{(117)}. Aqueles obrièrs d'amont, amb sas largas bragas de velós,
son fustièrs de gròs. Un ten una biga pichona e l'autre talha una cavilha amb sa
\VUgras{destral} \textsuperscript{(116)}. Aquel qu'es contra la
\VUgras{chiminèa} \textsuperscript{(41)} es lo teulejaire. Clava de listèls sus
las \VUgras{bigas} (*), e de \VUgras{lausas} \textsuperscript{(66)} suls listèls.
De còps i pausa de teules.

%%K t05-noto-1 noto
\VUnotes{143.2mm}{%
(*) \VUgras{Balk-o} = alemand \textit{Balken}, anglés \textit{joist}, francés
\textit{solive}, italian \textit{travicello}, rus \textit{balka}, espanhòl e
portugués \textit{viga}. Raiç tecnica pas encara adoptada, mas que se prepausa
aicí per traduire lo sens del mot francés \textit{solive}, qu'es pas un
\textit{trau} (francés \textit{poutre}) ni una biga pichona. Es un fust escairat
que sosten un sòl e un plafon.%
}

%%K t05-01-26 p
Lo teulat de l'estable es de \VUgras{teules} \textsuperscript{(55)}, aquel de
l'ostal es de \VUgras{lausa} \textsuperscript{(32)} e aquel de la galariá es de
\VUgras{zinc} \textsuperscript{(24)}.

%%K t05-01-27 p
\textsc{Joan}. --- Lo teulejaire fa tanben los teulats de zinc?

%%K t05-01-28 p
\textsc{Lo Sr. Blanc}. --- Non, aqueste es lo \VUgras{zinguièr}
\textsuperscript{(121)} o lo \VUgras{plombièr}, aquel qu'es a pausar una
\VUgras{luqueta} \textsuperscript{(38)} sul teulat de l'ostal. Pausa tanben las
\VUgras{andanas} \textsuperscript{(43)} e los \VUgras{davalaires}
\textsuperscript{(42)}. Solda lo zinc e lo lauton. Es el que fixèt lo
\VUgras{parafòldre} \textsuperscript{(40)} e la \VUgras{giroleta}
\textsuperscript{(39)} sul \VUgras{frontal} \textsuperscript{(36)}.

%%K t05-01-29 p
\textsc{Joan}. --- Alara l'ostal es acabat?

%%K t05-tit-5 sub
\VUsaut{3.0mm}
\VUpk{10.2pt}{40}{Los espleches.}
\VUsaut{3.0mm}

%%K t05-01-30 p
\textsc{Lo Sr. Blanc}. --- Pas del tot. Lo \VUgras{gipsièr}
\textsuperscript{(99)} auça amb de \VUgras{brica} \textsuperscript{(61)} las
paretas que la partisson dins las diferentas cambras. Revestís de \VUgras{gip}
\textsuperscript{(70)} los \VUgras{plafons} \textsuperscript{(26)} e las parets.
Ara es a far lo plafon de la \VUgras{galariá} \textsuperscript{(33)}. A, coma lo
maçon, una \VUgras{truèla} \textsuperscript{(100)} e una \VUgras{gaveta}
\textsuperscript{(101)}.

%%K t05-01-31 p
Puèi l'ebenista fa las pòrtas, las fenèstras, los \VUgras{parquets}
\textsuperscript{(16)} e los \VUgras{contravents} \textsuperscript{(19)}.

%%K t05-01-32 p
Ara trabalha dins la cort a son \VUgras{banc de fustièr} \textsuperscript{(90)}. A
una \VUgras{rèssa} \textsuperscript{(94)} per rèsser la \VUgras{fusta}
\textsuperscript{(69)}, un \VUgras{rabòt} \textsuperscript{(91)}, un
\VUgras{sarjant} \textsuperscript{(92)}, un \VUgras{cisèl}
\textsuperscript{(94 \textit{bis})}, un \VUgras{compàs} \textsuperscript{(95)} e
un \VUgras{maçòl de fusta} \textsuperscript{(95 \textit{bis})}.

%%K t05-01-33 p
\textsc{Joan}. --- A, mas es encara mai divertent d'èsser ebenista que maçon!
Oncle, me cromparàs un banc, una rèssa e un rabòt? Farai una cabana per Medor.

%%K t05-01-34 p
\textsc{L'Oncle}. --- Òc, filh meu.

%%K t05-01-35 p
\textsc{Joan}. --- Que soi content! E aquel qu'es contra la pòrta d'intrada, qual
es?

%%K t05-tit-6 sub
\VUsaut{3.0mm}
\VUpk{10.2pt}{40}{La pintura.}
\VUsaut{3.0mm}

%%K t05-01-36 p
\textsc{Lo Sr. Blanc}. --- Es un \VUgras{pintre} \textsuperscript{(102)}. Es sus
son escala amb son pòt de \VUgras{pintura} \textsuperscript{(103)} e son
\VUgras{brossa} \textsuperscript{(104)}. Pinta la fusta de la pòrta d'intrada per
que se conserve mai longtemps.

%%K t05-01-37 p
Es tanben el qu'empega lo papièr dins las cambras.

%%K t05-01-38 p
Lo \VUgras{tapissièr} \textsuperscript{(107)} qu'es sus una \VUgras{escala}
\textsuperscript{(106)}, sul \VUgras{balcon} \textsuperscript{(28)}, pausa una
\VUgras{marquisa} \textsuperscript{(29)}.

%%K t05-01-39 p
L'obrièr qu'es contra la granda fenèstra es lo \VUgras{veirièr}
\textsuperscript{(96)}. Fixa los \VUgras{veires} \textsuperscript{(18)} amb de
\VUgras{mastic} \textsuperscript{(73)}. Aquel autre obrièr, ajolhat contra la
pòrta de la cava, es un \VUgras{sarralhièr} \textsuperscript{(97)}. Es a pausar
una \VUgras{sarralha} \textsuperscript{(6)}. Sa \VUgras{lima}
\textsuperscript{(98)} es a son costat.

%%K t05-01-40 p
\textsc{Joan}. --- Es tanben sarralhièr aquel que tusta amb son \VUgras{martèl}
\textsuperscript{(84)} sus un fèrre?

%%K t05-01-41 p
\textsc{Lo Sr. Blanc}. --- Puslèu es un \VUgras{fabre} \textsuperscript{(125)}.
Fòrja de \VUgras{fèrres} \textsuperscript{(67)} per l'\VUgras{electricista}
\textsuperscript{(124)}, qu'es sul teulat de la \VUgras{remesa}
\textsuperscript{(51)}. A una \VUgras{fornal} \textsuperscript{(126)}, una
\VUgras{enclutge} \textsuperscript{(127)} e de \VUgras{tenalhas}
\textsuperscript{(128)}.

%%K t05-01-42 p
L'electricista pausa lo \VUgras{fial de coire} \textsuperscript{(44)} del telefòn.

%%K t05-tit-7 sub
\VUpk{10.2pt}{40}{La jardinariá.}
\VUsaut{3.0mm}

%%K t05-01-43 p
\textsc{L'Oncle}. --- Al fons de la \VUgras{cort} \textsuperscript{(45)}, contra
la \VUgras{reissa} \textsuperscript{(47)} e lo \VUgras{portal}
\textsuperscript{(48)}, veses lo \VUgras{jardinièr} \textsuperscript{(129)}. Es a
preparar lo pichon \VUgras{jardin} \textsuperscript{(49)}. A cavat la tèrra amb sa
\VUgras{bechada} \textsuperscript{(131)}, semena de granas e rastèla amb lo
\VUgras{rastèl} \textsuperscript{(130)}. Puèi, amb son \VUgras{arrosador}
\textsuperscript{(132)}, arrosarà las \VUgras{platabandas}
\textsuperscript{(150)} e las plantas creisseràn.

%%K t05-01-44 p
Lo \VUgras{varlet d'estable} \textsuperscript{(133)}, que ven de suènhar lo caval
e de li donar a manjar, neteja la \VUgras{veitura} \textsuperscript{(52)} amb una
esponga, una \VUgras{bomba a man} \textsuperscript{(134)} e un selh d'aiga. Puèi
la butarà dins la remesa e barrarà la pòrta amb lo gròs \VUgras{ferrolh}
\textsuperscript{(53)}.

%%K t05-01-45 p
Ara siás content, supausi; as agut pro d'explicacions.

%%K t05-01-46 p
\textsc{Joan}. --- Òc, oncle, mas m'agradariá de veire l'interior de l'ostal.

%%K t05-01-47 p
\textsc{L'Oncle}. --- Aquò serà deman. Lo sénher Roge t'explicarà lo plan e te
dirà lo nom de cada cambra.

%%K t05-tit-8 sub
\VUsaut{3.0mm}
\VUpk{10.2pt}{40}{La distribucion interiora de l'ostal.}
\VUsaut{2.4mm}

%%K t05-apar-2 apar
{\centering\textit{(Vejatz lo plan.)}\par}
\VUsaut{2.4mm}

%%K t05-01-48 p
\textsc{L'Arquitècte}. --- Vèni, Joan, te farai visitar las diferentas partidas de
l'ostal.

%%K t05-01-49 p
Dintrem al \VUgras{ras-de-caussada} \textsuperscript{(7)}. Vaicí lo boton de la
\VUgras{sonalha} \textsuperscript{(11)}. Pujam los tres \VUgras{escalons}
\textsuperscript{(14)} del \VUgras{perron} \textsuperscript{(9)}, passam lo
\VUgras{lindal} \textsuperscript{(10)} de la \VUgras{pòrta d'intrada}
\textsuperscript{(8)}, e vaicí que sèm al pè de l'\VUgras{escalièr}
\textsuperscript{(13)}.

%%K t05-01-50 p
\textsc{Joan}. --- Aquò's lo corredor.

%%K t05-01-51 p
\textsc{L'Arquitècte}. --- Non, aquò's lo \VUgras{vestibul}
\textsuperscript{(12)}. Lo \VUgras{corredor} \textsuperscript{(a)} es al fons. A
drecha, vaicí lo \VUgras{salon} \textsuperscript{(b)} e la \VUgras{sala de
manjar} \textsuperscript{(c)}; davant la sala de manjar i a la \VUgras{cosina}
\textsuperscript{(d)} e l'\VUgras{ofici} \textsuperscript{(e)}; puèi, davant, lo
\VUgras{cabinet de trabalh} \textsuperscript{(f)}. La \VUgras{galariá}
\textsuperscript{(23)}, amb sa \VUgras{pòrta veirada} \textsuperscript{(25)}, es
contra lo salon.

%%K t05-01-52 p
Ara pujarem l'escalièr. Agarra-te a la \VUgras{rampa} \textsuperscript{(15)} e
compta los escalons. Ne son quatòrze. Vaicí lo \VUgras{palièr}
\textsuperscript{(m)}: sèm al primièr \VUgras{estatge} \textsuperscript{(27)}.

%%K t05-tit-9 sub
\VUsaut{3.0mm}
\VUpk{10.2pt}{40}{Lo primièr estatge.}
\VUsaut{3.0mm}

%%K t05-01-53 p
Davant l'escalièr i a los \VUgras{comuns} \textsuperscript{(20)} e lo
\VUgras{cabinet de toaleta} \textsuperscript{(j)}. A drecha, una bèla
\VUgras{cambra} \textsuperscript{(g)} amb doas fenèstras sus la \VUgras{façada}
\textsuperscript{(1)} de l'ostal. A esquèrra i a la \VUgras{sala de banh}
\textsuperscript{(1)} e la \VUgras{cambra d'ostes} \textsuperscript{(i)} amb una
granda \VUgras{alcòva} \textsuperscript{(h)}. Contra la cambra i a la
\VUgras{cambra dels enfants} \textsuperscript{(k)}. Dona sus una larga
\VUgras{terrassa} \textsuperscript{(21)}. Jos aquela terrassa i a d'autres comuns
\textsuperscript{(20)} e, contra la paret, vaicí la \VUgras{campana}
\textsuperscript{(22)} que sona lo sopar.

%%K t05-01-54 p
\textsc{Joan}. --- Pugem al \VUgras{segond estatge}.

%%K t05-01-55 p
\textsc{L'Arquitècte}. --- I a pas de segond estatge. Jos lo teulat i a pas que de
\VUgras{mansardas} \textsuperscript{(33)} e las trèvas \textsuperscript{(34)}.
Avèm acabat la visita, podèm davalar.

%%K t05-01-56 p
\textsc{Joan}. --- Mercés, sénher, es estat fòrça interessant. Vau contar a mon
paire tot çò qu'ai vist.
\VUsaut{4.0mm}
\VUfilet{22mm}
